\documentclass{article}

\usepackage[preprint,nonatbib]{neurips_2020}
\makeatletter
\renewcommand{\@noticestring}{}
\makeatother

\usepackage[utf8]{inputenc}
\usepackage[T1]{fontenc}
\usepackage{hyperref}
\usepackage{url}
\usepackage{booktabs}
\usepackage{amsfonts}
\usepackage{amsmath}
\usepackage{nicefrac}
\usepackage{microtype}
\usepackage{graphicx}
\usepackage{placeins}

\title{Vector-Q Transformer SAC for Autonomous Racing in Assetto Corsa}

\author{%
  Chandan Gowda Keelara Shivanna \\
  Khoury College of Computer Sciences \\
  Northeastern University \\
  Boston, MA 02115 \\
  \texttt{gowdakeelarashivan.c@northeastern.edu}
}

\begin{document}

\maketitle

\begin{abstract}
  Soft Actor-Critic (SAC) has become a default choice for continuous-control
  reinforcement learning, but its scalar $Q$-function and scalar
  temperature $\alpha$ implicitly assume that action channels share a
  common credit signal. In autonomous racing this assumption is strained:
  steering, throttle, and brake are physically coupled but their effects
  on progress and lap-time are distinct, so a joint $Q$ blurs
  channel-specific credit. We present \emph{Vector-Q Transformer SAC}, a
  single-agent SAC variant that (i) decomposes the critic into
  per-action-channel $Q$ heads trained from a per-channel vector reward,
  (ii) assigns each channel its own auto-tuned temperature $\alpha$ with
  independently configurable target entropy, (iii) replaces the MLP
  encoder with a $4$-layer Transformer over a $75$-frame (3\,s) window,
  and (iv) seeds a $6$-channel stratified replay buffer with a
  behavior-cloning dataset of human-driven demonstrations. On Monza
  with a Mazda Miata in Assetto Corsa our method
  completes a full lap on $80\%$ of evaluation episodes, whereas a
  tuned MLP SAC baseline never completes a lap, reaching at best
  $27\%$ of the track.
  Per-frame reward on positive frames is comparable between the two,
  indicating that the main gain is in the \emph{persistence} of good
  driving rather than its peak quality.
\end{abstract}

\section{Introduction}
\label{sec:intro}

Autonomous racing is a stringent continuous-control benchmark: a policy
must output precise steering, throttle, and brake at $25$\,Hz, reason
over long horizons in which dense per-frame rewards must accumulate
into lap-time, and tolerate sensor noise and latency from a commercial
simulator. The three action channels are physically coupled (braking
off-line amplifies lateral error; throttling mid-corner induces
understeer), yet causally distinct: the right brake action may coexist
with a wrong throttle action, and a scalar joint reward obscures this.

A natural baseline is Soft Actor-Critic (SAC)~[1] with automatic
temperature tuning~[2]: a scalar $Q$ is trained on a scalar reward and
a single $\alpha$ regularises the joint policy entropy. We observe
that this formulation is limiting on our task. A tuned MLP SAC
baseline trained on Assetto Corsa never completed a lap; its best
episode covered only $27\%$ of Monza.

We propose \textbf{Vector-Q Transformer SAC}, a single-agent SAC
variant with four architectural changes aimed at per-channel credit
assignment and temporal context:
\begin{enumerate}
  \item A per-channel critic $[Q^{\mathrm{s}}, Q^{\mathrm{th}}, Q^{\mathrm{br}}]$
  trained from a per-channel vector reward, adapting cooperative-MARL
  value decomposition~[5,6] to single-agent continuous control.
  \item A per-channel temperature $\alpha_c$ with an independent
  target entropy $H^c_{\text{target}}$, letting each channel find its
  own exploration/exploitation equilibrium.
  \item A $4$-layer, $4$-head, $d_{\text{model}}{=}256$ Pre-LN
  Transformer observation encoder over a $75$-frame ($3$\,s) window,
  giving the policy explicit corner-approach context.
  \item A $6$-channel stratified FIFO replay ($10^5$ per bin) seeded
  with $18{,}497$ behavior-cloning (BC) demonstration windows,
  sharpening the critic toward the expert racing line.
\end{enumerate}

Our method completes a full lap on $80\%$ of evaluation episodes,
drives $6.2\times$ further per episode than the baseline at
$1.4\times$ the mean speed and $2.4\times$ the per-frame reward.

\paragraph{Scope note.}
An earlier proposal outlined TD3~[3] with an MLP policy at $60$\,Hz;
early experiments motivated a switch to SAC (stabler under our
reward) and to a Transformer encoder (temporal context was the
dominant bottleneck), and we standardised on the $25$\,Hz control
rate the plugin supports without packet loss.

\paragraph{Open-source code disclosure.}
Our training, Vector-$Q$ architecture, per-channel $\alpha$ logic,
stratified replay buffer, behavior-cloning injection, and reward
pipeline are original contributions. We use the open-source
\texttt{assetto\_corsa\_gym}~[14] plugin as the data bridge to
Assetto Corsa.

\section{Related Work}
\label{sec:related}

\paragraph{Soft Actor-Critic and auto-tuned temperature.}
SAC~[1] is an off-policy maximum-entropy actor-critic method; its
objective balances reward against entropy via a scalar $\alpha$. The
automatic-tuning extension~[2] drives $\mathbb{E}[\log\pi(a|s)]$
toward a scalar target entropy $H_\text{target}$, typically
$-|\mathcal A|$. Both works use a single $\alpha$ and a scalar
$Q(s,a)$ on a joint action. Our method retains the SAC loss structure
but relaxes both along the action dimension.

\paragraph{Value decomposition.}
VDN~[5] and QMIX~[6] decompose a cooperative team $Q$ into per-agent
summands (or a monotonic mixture) so each agent receives a direct
gradient from its return share, reporting sample-efficiency gains
attributable to per-agent credit. In single-agent RL, the Hybrid
Reward Architecture~[11] factorises a scalar reward into components
and trains one $Q$-head per component. Our Vector-$Q$ critic is
closest in spirit to HRA but decomposes per \emph{action channel}
rather than per reward component, trained from a per-channel vector
reward. Per-channel $\alpha$ does not alter the critic's gradient
structure and is an entropy-schedule expressiveness mechanism
more than a sample-efficiency mechanism.

\paragraph{Multi-agent framings.}
Our decomposition is inspired by MARL but is not multi-agent: we use
a shared encoder, a joint policy, and jointly sampled actions.
MPO~[4]'s per-state KL constraint is conceptually related to
per-channel entropy targets but bounds policy step-size rather than
decomposing credit.

\paragraph{Imitation and advantage-weighted methods.}
DQfD~[13] injects demonstrations into an off-policy replay buffer for
$Q$-learning; we follow the same idea adapted to SAC and a
per-channel stratified buffer. DAgger~[9] motivates correcting
covariate shift via interactive expert mixing, which we do not
employ. AWAC~[10] advantage-weights demonstration actions in the
policy update; we rely on the off-policy SAC gradient to absorb the
injected demonstrations through the critic rather than re-weighting
the policy loss directly.

\paragraph{Transformers in RL and racing-game RL.}
Transformers have been used both as full sequence models, notably
Decision Transformer~[8], which reframes RL as return-conditioned
supervised learning, and as drop-in observation encoders. Our use is
strictly the latter: the Transformer replaces SAC's MLP encoder
without changing the off-policy loss, so the entropy-regularised
actor-critic guarantees are preserved. Fuchs et al.~[7] trained a
scalar-$Q$ SAC policy to super-human performance on Gran Turismo
Sport using large-scale parallel simulation; we share their
observation that reward shaping and a stable replay buffer are
decisive for lap completion but operate at much smaller scale and
on a commercial simulator with sensor latency, and we additionally
decompose both the critic and the temperature along the action
dimension.

\section{Methods}
\label{sec:methods}

\begin{figure}[t]
  \centering
  \includegraphics[width=0.88\linewidth]{figures/architecture.png}
  \caption{\textbf{Vector-Q Transformer SAC architecture.} Each frame
  emits a $50$-dim observation token; $75$ tokens stack into a
  $3$-second window. Two independent Transformer encoders (policy and
  critic, same architecture, different weights) produce a $256$-dim
  embedding. The policy head (left) outputs a diagonal-Gaussian
  tanh-squashed action with per-channel auto-tuned temperatures
  $\alpha_\mathrm{s}, \alpha_\mathrm{th}, \alpha_\mathrm{br}$ each
  targeting its own entropy. The twin critic head (right) concatenates
  the embedding with the action and produces a per-channel vector
  $\mathbf{Q}(s,a)\in\mathbb{R}^3$ (clipped double-$Q$).}
  \label{fig:arch}
\end{figure}

\subsection{Problem setting}

The simulator runs at $25$\,Hz. Each tick emits a $62$-dimensional
state vector; we retain the first $50$ dimensions as a per-frame
observation token $o_t \in \mathbb{R}^{50}$ (car-state channels, ray
casts to the track boundary, OOT flag, curvature look-ahead,
past-action history), and stack $W{=}75$ tokens to form
$O_t \in \mathbb{R}^{W\times 50}$. The action
$a_t = (a^\mathrm{s}, a^\mathrm{th}, a^\mathrm{br}) \in [-1, 1]^3$ maps
to Assetto Corsa's steering, throttle, and brake controls. A scalar
reward $r_t \in [-1, 1]$ is a weighted sum of five normalised
components (progress, speed-tracking, absolute gap to racing line,
action smoothness, yaw alignment) with crash frames short-circuited
to $-1$. These components are not all aligned across action
channels: for instance, applying brake when the car is moving
slowly deserves a penalty while applying throttle in the same
state deserves a reward, and a scalar sum averages these two
opposing signals into one number that largely cancels out. We
therefore decompose the reward into a per-channel vector
$\mathbf{r}_t = (r^\mathrm{s}_t, r^\mathrm{th}_t, r^\mathrm{br}_t)$,
where each channel reweights the same underlying components so
that opposing signals on different action dimensions are preserved
rather than collapsed.

\subsection{Baseline SAC}

SAC trains a policy $\pi_\phi$ and clipped-double $Q$
$Q_\theta = \min_{i=1,2} Q_{\theta_i}$ by minimising
\begin{align}
\mathcal{L}_Q(\theta_i) &= \mathbb{E}
\!\left[\left(Q_{\theta_i}(s,a) - r - \gamma(1{-}d)\!\left[Q_{\bar\theta}(s', a') -
  \alpha \log \pi_\phi(a'|s')\right]\right)^2\right],
\label{eq:sac-q}\\
\mathcal{L}_\pi(\phi) &= \mathbb{E}\bigl[\alpha \log \pi_\phi(a|s) - Q_\theta(s, a)\bigr],
\;\;
\mathcal{L}_\alpha(\alpha) = -\mathbb{E}\bigl[\alpha(\log\pi_\phi(a|s) + H_\text{target})\bigr].
\label{eq:sac-pi}
\end{align}
The scalar $\alpha$ applies a uniform entropy penalty across all
action dimensions, and the scalar $Q$ gradient flows uniformly into
the actor's output neurons.

\subsection{Vector-$Q$ critic decomposition}
\label{sec:vector-q}

\emph{Why.} A scalar $Q$ receives a single gradient from a single
reward and distributes it uniformly across the three output action
neurons. When the per-channel reward contributions point in opposing
directions on the same frame, a scalar $Q$ cannot represent "the
brake head was right, the throttle head was wrong" — both heads
receive the same blended signal. Decomposing the critic along the
action dimension gives each channel its own gradient from its own
reward.

We replace the scalar critic head with a vector
$\mathbf{Q}_\theta(s,a) = [Q^\mathrm{s}, Q^\mathrm{th}, Q^\mathrm{br}] \in \mathbb{R}^3$
trained against per-channel Bellman targets:
\begin{equation}
\mathcal{L}_Q^{\mathrm{vec}}(\theta_i) = \mathbb{E}
\!\left[\sum_{c}
  \left(Q^c_{\theta_i}(s,a) - r^c - \gamma(1{-}d)\!\left[Q^c_{\bar\theta}(s', a') - \alpha_c \log \pi^c_\phi(a'|s')\right]\right)^2\right],
\label{eq:vec-q}
\end{equation}
with policy loss
$\mathcal{L}_\pi^{\mathrm{vec}}(\phi) = \mathbb{E}\bigl[\sum_c \alpha_c \log \pi^c_\phi(a^c|s) - Q^c_\theta(s, a)\bigr]$.
Each critic output neuron now receives a gradient from a
channel-specific reward, so a correctly braking frame with
simultaneously wrong throttle no longer penalises the brake $Q$-head.
The construction adapts the VDN/QMIX credit-assignment intuition~[5,6]
to single-agent continuous control.

\subsection{Per-channel temperature}
\label{sec:per-channel-alpha}

\emph{Why.} The three action channels have qualitatively different
exploration needs: steering benefits from residual stochastic
correction, while throttle and brake are closer to bimodal and
benefit from decisive, low-entropy decisions. A single scalar
$\alpha$ forces all three channels to share one exploration level
and must settle at a compromise that is wrong for at least one of
them. Auto-tuning $\alpha$ per channel removes this constraint.

The scalar $\alpha$ is replaced by a per-channel vector
$\boldsymbol{\alpha} = (\alpha_\mathrm{s}, \alpha_\mathrm{th}, \alpha_\mathrm{br})$
with $\alpha_c = \exp(\log\alpha_c)$ clamped to
$\log\alpha_c \in [-5, 0]$. The auto-tuning loss is channel-wise,
\begin{equation}
\mathcal{L}_\alpha(\alpha_c) = -\mathbb{E}\bigl[\alpha_c(\log\pi^c_\phi(a^c|s) + H^c_\text{target})\bigr],
\label{eq:per-channel-alpha}
\end{equation}
with targets $H_\text{target} = (-2, -3, -3)$: looser on steering
(benefits from stochastic correction) and tighter on throttle and
brake (mostly bimodal, benefit from determinism).

\subsection{Transformer observation encoder}
\label{sec:encoder}

\emph{Why.} Corner-entry decisions are not instantaneous functions
of the current frame: the correct brake point depends on how speed
has been evolving over the last few seconds, and the correct
steering angle depends on the yaw trajectory that led into the
corner. An MLP on a single-frame observation has no access to this
history and must rediscover it through the value function. We
instead give the policy and critic explicit temporal context with
a Transformer encoder over a $3$-second window, turning the
corner-approach history into an input feature rather than a
credit-assignment problem.

The full architecture is shown in Figure~\ref{fig:arch}. The encoder
applies a linear projection to $d{=}256$, a learned
positional embedding over $75$ positions, four Pre-LN
Transformer~[12] layers with $4$ heads and $1024$-wide feed-forward,
and mean-pooling
over the sequence to yield $h_t \in \mathbb{R}^{256}$. The policy and
critic encoders are independent copies to prevent joint-embedding
collapse. The policy head is a two-layer MLP
$256\!\to\!256\!\to\!6$ producing a diagonal-Gaussian tanh-squashed
action; each critic head is $(256{+}3)\!\to\!256\!\to\!3$.

\subsection{Stratified replay with BC demonstration seeding}
\label{sec:replay}

\emph{Why.} Most frames on a lap are straight-line, high-throttle,
high-reward, so a uniform FIFO buffer is dominated by easy positive
states and the critic rarely sees informative failure states. And
at the start of training the policy has no idea what a competent
racing line looks like, so the critic has no expert trajectories to
score. We address both with the replay design: stratification
enforces per-channel positive/negative balance inside every
minibatch, and a one-time behavior-cloning seed injects expert
states into the positive sub-buffers before the first gradient
step.

Transitions are stored in a six-channel stratified FIFO buffer: each
action channel has a positive and negative sub-buffer ($r^c{>}0$ vs
$r^c{\le}0$) of capacity $10^5$ windows, giving $6{\times}10^5$
total. Minibatches sample uniformly from each sub-buffer to counteract
the positive-reward imbalance on straights. Before training we seed
the positive sub-buffers with $18{,}497$ pre-collected behavior-cloning
demonstration windows; this one-time injection sharpens the critic on
the expert racing line without modifying the policy loss. Optimisation
uses Adam at $3{\times}10^{-4}$, $\gamma{=}0.992$, Polyak
$\tau{=}0.005$, batch size $256$.

\paragraph{Parameter budget.}
The full model has $6.6$M parameters, $\approx 97\%$ concentrated in
the two independent Transformer encoders; the policy and twin-critic
heads add only $\approx 200$k. A matched MLP SAC with hidden
$[256,256,256]$ on a single-frame observation has $0.44$M parameters,
so our encoder accounts for essentially all of the additional capacity
(a $15\times$ increase), purchased for the $3$-second temporal
context. The per-channel vector critic and per-channel $\alpha$ are
near-free on top of the encoders, adding three scalar $\alpha_c$
parameters and widening each critic output from $1$ to $3$.

\section{Experiments}
\label{sec:experiments}

\subsection{Testbed and research questions}

We evaluate on Monza with a Mazda Miata. Each episode begins with
a $25$-second scripted warm-up that randomises the start position
along the racing line, followed by a $150$--$250$-step
straight/full-throttle warm-up, then the learned policy takes
control. Episodes terminate on OOT, stationary-car ($<0.5$\,m in
$60$ frames), or a max-frame cap. We ask four questions:
\textbf{Q1} does our policy reliably complete laps?
\textbf{Q2} does the BC injection measurably shift the critic?
\textbf{Q3} does per-channel target-entropy tightening yield
channel-specific decisiveness?
\textbf{Q4} where does the policy still fail?

\subsection{Main result}

\begin{figure}[!h]
  \centering
  \includegraphics[width=\linewidth]{figures/trajectory.png}
  \caption{\textbf{Speed and brake-zone map for a Vector-Q
  Transformer SAC full-lap evaluation episode on Monza.}
  (a) Trajectory coloured by per-frame speed; the policy reaches
  $\sim\!46$\,m/s on the straights and slows to $13$--$20$\,m/s at
  each chicane. (b) Trajectory coloured by per-frame brake action;
  braking is sharply concentrated at five to six distinct zones per
  lap, each aligned with a named Monza corner.}
  \label{fig:traj}
\end{figure}

Table~\ref{tab:main} compares the baseline against our variant. The
baseline is an MLP SAC with hidden $[256, 256, 256]$, scalar $Q$,
scalar $\alpha$, target entropy $-3$.

\begin{table}[h]
\centering
\footnotesize
\caption{Baseline MLP SAC vs Vector-Q Transformer SAC on Monza.}
\label{tab:main}
\begin{tabular}{@{}lccc@{}}
\toprule
metric & Baseline MLP SAC & Vector-Q Transformer SAC & $\Delta$ \\
\midrule
Lap completion rate                    & $0\%$        & $\mathbf{80\%}$         & $+80$\,pp \\
Mean distance / episode (m)            & $619$        & $\mathbf{3{,}862}$      & $6.2\times$ \\
Best episode distance (m)              & $1{,}539$    & $\mathbf{5{,}774}$      & $3.8\times$ \\
Speed mean (m/s)                       & $19.9$       & $\mathbf{28.8}$         & $1.4\times$ \\
Speed max (m/s)                        & $43.1$       & $44.1$                  & $\approx$ \\
Mean $|$gap$|$ (m)                     & $0.48$       & $0.52$                  & $\approx$ \\
Reward / frame                         & $+0.27$      & $\mathbf{+0.64}$        & $2.4\times$ \\
Positive-reward frames (\%)            & $97$         & $97.9$                  & $\approx$ \\
\bottomrule
\end{tabular}
\end{table}

Both policies track the racing line with comparable per-frame
precision ($\approx 0.5$\,m lateral error) and reach similar top
speeds on the frames they actually drive; the difference is in
\emph{persistence}. Our variant sustains good driving for much
longer, yielding $6.2\times$ the distance per episode and completing
full laps on the majority of evaluation runs. Per-frame reward is
$2.4\times$
higher at a comparable positive-frame fraction, indicating the
policy is selecting higher-value actions on those frames.
Qualitatively, a full-lap evaluation produces five to six distinct
brake zones aligned with Monza's named corners; speed peaks near
$46$\,m/s on the straights and drops to $13$--$20$\,m/s at each
chicane, and brake application is sharply binary rather than
modulated (Figure~\ref{fig:traj}). This answers \textbf{Q1}
qualitatively: the policy drives a full lap with a speed and brake
profile consistent with a reasonable racing line.

\subsection{Ablation and analysis}
\label{sec:ablation}

For each of the four mechanisms we present a targeted diagnostic
designed to isolate one behavioural claim about that mechanism and
to rule out a specific failure mode it could exhibit.

\paragraph{BC injection (Q2).}
After the one-time injection and several training epochs, per-channel
$Q$ on positive-sample states increases by $11$--$15\%$ on all three
channels while negative-sample $Q$ moves by less than $3\%$.
Evaluated directly on the demonstration set, demo actions had
$Q$-values statistically indistinguishable from the policy's own
before injection ($\Delta \approx 0.1\%$ of $Q$ magnitude, effectively
flat gradient); after injection they are pulled measurably above
policy actions on all channels, giving the actor a usable gradient
toward the expert line.

\paragraph{Per-channel $\alpha$ reaches distinct equilibria (Q3).}
Figure~\ref{fig:alpha} plots the three per-channel temperatures
$\alpha_c$ and policy entropies $H_c$ over $\sim\!5000$ gradient
steps of end-of-training fine-tuning. The three channels reach
clearly separated $\alpha$ equilibria
($\alpha_\mathrm{brake} \gg \alpha_\mathrm{throttle} \gg \alpha_\mathrm{steer}$),
with $\alpha_\mathrm{steer}$ pinned near the $e^{-5}$ clamp floor and
$\alpha_\mathrm{brake}$ settling roughly $3.5\times$ higher. The
corresponding policy entropies equilibrate to visibly different
values on each channel as well. A single scalar $\alpha$ cannot
simultaneously satisfy this per-channel spread: it would collapse to
a compromise value that deviates from at least one channel's natural
equilibrium by a multiplicative factor of roughly $3.5\times$. In
the same training budget, our per-channel auto-tuner reaches a
stable multi-channel equilibrium from the first gradient step,
whereas a scalar auto-tuner would need additional steps to stabilise
whichever channels its compromise misaligned. The per-channel
mechanism is therefore doing work that its scalar counterpart
cannot, and is consistent with improved training efficiency under
equal compute.

\begin{figure}[!h]
  \centering
  \includegraphics[width=0.95\linewidth]{figures/per_channel_alpha.png}
  \caption{\textbf{Per-channel $\alpha$ auto-tuning reaches distinct
  equilibria on each action channel.}
  (a) $\alpha_c$ over training; the three channels separate into
  a clear ordering
  $\alpha_\mathrm{brake} \gg \alpha_\mathrm{throttle} \gg \alpha_\mathrm{steer}$,
  with $\alpha_\mathrm{steer}$ pinned near the clamp floor.
  (b) Corresponding policy entropies $H_c$. The differentiation
  justifies per-channel temperature over a single scalar $\alpha$.}
  \label{fig:alpha}
\end{figure}

\paragraph{Per-channel $Q$-discrimination.}
Figure~\ref{fig:qdist} shows the distribution of learned $Q$-values
on each of the three critic heads over held-out buffer samples.
All three channels exhibit a clear bimodal structure: a low-value
mode near $Q\!\approx\!0$ and a high-value mode near $Q\!\approx\!90$.
The bimodality reflects the positive/negative stratification of the
training buffer and confirms that every critic head has learned a
meaningful score separating favourable from unfavourable states.
Splitting $Q$ on each head by the sign of the per-channel reward
sharpens this into a quantitative statement: the
positive-over-negative $Q$ ratio is $\approx 4.7\times$ on steer,
$4.8\times$ on throttle, $4.7\times$ on brake (see
Appendix~\ref{app:qpn}). The near-identical ratios across channels
are evidence that the decomposition is \emph{balanced}: no single
channel dominates the critic's discriminative power. A naive
vector-$Q$ can collapse into a degenerate regime in which one head
captures nearly all of the learning signal and the others remain
close to random (a failure mode in which vector-$Q$ behaves as a
scalar-$Q$ in disguise); our observed $4.7$--$4.8\times$ ratios on
all three channels show that this collapse did not occur.

\begin{figure}[!h]
  \centering
  \includegraphics[width=\linewidth]{figures/q_distribution.png}
  \caption{\textbf{Per-channel $Q$-value distributions} on held-out
  buffer samples, one histogram per critic head. Each channel shows
  a pronounced bimodality, separating good from bad states.}
  \label{fig:qdist}
\end{figure}

\paragraph{Residual failure (Q4).}
The one recurrent failure is the Variante della Roggia chicane at
race pace. On terminated episodes the car is $1.5$--$2$\,m off-line
at the braking point and the brake pedal reaches only $\approx 0.7$
of its observed maximum, whereas Roggia requires a full brake. The
corner is under-represented in the demonstration set, and the policy
internalised the softer brake.

\FloatBarrier
\section{Conclusion}
\label{sec:conclusion}

\paragraph{Summary.}
Vector-Q Transformer SAC combines per-channel critic decomposition,
per-channel auto-tuned temperature, a $3$-second Transformer context
window, and BC-demonstration buffer seeding. On Monza it completes a
lap on $80\%$ of evaluation episodes, drives $6.2\times$ further
per episode than the best MLP SAC baseline
available for this environment, and matches its per-frame lateral
precision.

\paragraph{Limitations.}
The policy is not yet perfect at the Variante della Roggia chicane
when approached at race pace, behaving well elsewhere but
under-braking here; we expect further training with a more diverse
set of examples at that corner to close the gap.

\paragraph{Reward-sensitive automatic $\alpha$ tuning (primary future work).}
Our most promising direction is making $\alpha_c$'s target entropy
itself state-dependent. In standard SAC and in our per-channel
variant, $\alpha_c$ is auto-tuned against a \emph{fixed}
$H^c_\text{target}$ that does not depend on the state. In practice
the right amount of exploration clearly depends on where the agent
is: on states where the current policy is already earning high
reward $\alpha_c$ should collapse toward its floor so the policy
becomes crisply deterministic, and on states where the policy is
struggling (low per-frame reward, large negative TD-error, or high
critic disagreement) $\alpha_c$ should rise and the policy should
re-explore the local action space. The SAC derivation admits this
extension directly: replace $H^c_\text{target}$ in
Eq.~\eqref{eq:per-channel-alpha} with a state-conditioned target
$H^c_\text{target}(s) = H_0 - \kappa\tanh\bigl(A^c(s,a)/\tau\bigr)$,
where $A^c{=}Q^c - V^c$ is the local per-channel advantage and
$H_0, \kappa, \tau$ are hyper-parameters. This tightens entropy on
high-advantage states and loosens it on low-advantage states, turning
the auto-tuning dual-update into a state-specific control loop rather
than a single global knob. The proposal is conceptually related to
MPO~[4]'s per-state KL constraint but operates on the entropy target
rather than the KL bound, and it preserves SAC's scalability. Full
specification, convergence analysis, and empirical validation are
left to future work.

\paragraph{Additional directions.}
(i) a corner-specific curriculum (targeted re-collection at Roggia);
(ii) multi-track generalisation to test whether the Transformer
context window and Vector-$Q$ gains survive track-level distribution
shift;
(iii) critic ensembles composed with the per-channel $Q$ architecture,
providing an uncertainty signal the reward-sensitive target above
can consume in place of the advantage.

\section*{References}

\small

[1] Haarnoja, T., Zhou, A., Abbeel, P., \& Levine, S.\ (2018) Soft
Actor-Critic: Off-Policy Maximum Entropy Deep Reinforcement Learning
with a Stochastic Actor. In \emph{Proc.\ ICML}.

[2] Haarnoja, T., Zhou, A., Hartikainen, K., Tucker, G., et al.\
(2018) Soft Actor-Critic Algorithms and Applications.
\emph{arXiv:1812.05905}.

[3] Fujimoto, S., van Hoof, H., \& Meger, D.\ (2018) Addressing
Function Approximation Error in Actor-Critic Methods (TD3). In
\emph{Proc.\ ICML}.

[4] Abdolmaleki, A., Springenberg, J.\ T., Tassa, Y., Munos, R.,
Heess, N., \& Riedmiller, M.\ (2018) Maximum a Posteriori Policy
Optimisation. In \emph{Proc.\ ICLR}.

[5] Sunehag, P., Lever, G., Gruslys, A., Czarnecki, W.\ M., et al.\
(2017) Value-Decomposition Networks for Cooperative Multi-Agent
Learning. \emph{arXiv:1706.05296}.

[6] Rashid, T., Samvelyan, M., Schroeder, C., Farquhar, G., Foerster,
J., \& Whiteson, S.\ (2018) QMIX: Monotonic Value Function
Factorisation for Deep Multi-Agent Reinforcement Learning. In
\emph{Proc.\ ICML}.

[7] Fuchs, F., Song, Y., Kaufmann, E., Scaramuzza, D., \& Duerr, P.\
(2021) Super-Human Performance in Gran Turismo Sport Using Deep
Reinforcement Learning. \emph{IEEE RA-L} \textbf{6}(3): 4257--4264.

[8] Chen, L., Lu, K., Rajeswaran, A., Lee, K., Grover, A., Laskin, M.,
Abbeel, P., Srinivas, A., \& Mordatch, I.\ (2021) Decision
Transformer: Reinforcement Learning via Sequence Modeling. In
\emph{Proc.\ NeurIPS}.

[9] Ross, S., Gordon, G.\ J., \& Bagnell, J.\ A.\ (2011) A Reduction
of Imitation Learning and Structured Prediction to No-Regret Online
Learning. In \emph{Proc.\ AISTATS}.

[10] Nair, A., Gupta, A., Dalal, M., \& Levine, S.\ (2020) AWAC:
Accelerating Online Reinforcement Learning with Offline Datasets.
\emph{arXiv:2006.09359}.

[11] van Seijen, H., Fatemi, M., Romoff, J., Laroche, R., Barnes, T.,
\& Tsang, J.\ (2017) Hybrid Reward Architecture for Reinforcement
Learning. In \emph{Proc.\ NeurIPS}.

[12] Vaswani, A., Shazeer, N., Parmar, N., Uszkoreit, J., Jones, L.,
Gomez, A.\ N., Kaiser, \L., \& Polosukhin, I.\ (2017) Attention Is
All You Need. In \emph{Proc.\ NeurIPS}.

[13] Hester, T., Vecerik, M., Pietquin, O., Lanctot, M., Schaul, T.,
Piot, B., Horgan, D., Quan, J., Sendonaris, A., Osband, I., et al.\
(2018) Deep Q-Learning from Demonstrations. In \emph{Proc.\ AAAI}.

[14] Remonda, A., Hansen, E., Raji, A., Graule, S., Bodinier, Q.,
et al.\ (2024) Assetto Corsa Gym: a reinforcement learning
environment for racing research.
\url{https://github.com/dasGringuen/assetto_corsa_gym}.

\normalsize

\clearpage
\appendix

\section{Additional critic diagnostics}
\label{app:critic}

\subsection{Per-channel $Q$ on positive vs negative sub-buffers}
\label{app:qpn}

Figure~\ref{fig:qpn} complements the per-channel $Q$ histograms in
Figure~\ref{fig:qdist} by splitting each head by the sign of the
per-channel reward. All three channels show clean separation with
$\mathrm{pos}/\mathrm{neg}$ ratios of $\approx 4.7$--$4.8\times$,
confirming that each $Q$-head has learned to score its own action
dimension and that no channel has collapsed.

\begin{figure}[h]
  \centering
  \includegraphics[width=\linewidth]{figures/q_pos_vs_neg.png}
  \caption{\textbf{Per-channel $Q$-value, positive vs negative
  buffer.} Box plots of the learned $Q$-values when the sampled
  transition carries a positive versus a negative per-channel
  reward.}
  \label{fig:qpn}
\end{figure}

\subsection{Per-channel $Q$ vs per-frame reward}
\label{app:q-vs-reward}

Figure~\ref{fig:qvr} shows the per-channel $Q$-value plotted against
the actual per-frame reward sampled from the training buffer, one
scatter per critic head. On all three channels the points cluster in
two visibly distinct regions (low reward aligned with low $Q$, high
reward aligned with high $Q$), and the upper cluster concentrates
near $Q \!\approx\! 90$--$100$. This is the finer-grained view of
the bimodality reported in Figure~\ref{fig:qdist}, and confirms that
the critic's learned value is positively correlated with the
channel-specific reward it was trained on.

\begin{figure}[h]
  \centering
  \includegraphics[width=\linewidth]{figures/q_vs_reward.png}
  \caption{\textbf{Per-channel $Q$-value versus actual per-frame
  reward}, one scatter per critic head. Points cluster into the two
  regimes the stratified buffer was designed to represent, and the
  high-reward high-$Q$ cluster is tight on all three channels,
  consistent with a well-trained per-channel critic.}
  \label{fig:qvr}
\end{figure}

\end{document}
